\section{Chương~\ref{sec:acet}. Thêm \nt{ACET}}

Ba tác động của acetylcholine đã được đo trên lát cắt và trong não sống, xung đột giữa ghi và đọc được chỉ ra trên mô hình của sách, còn việc \nt{ACET} làm đường cho phép ghi và báo sự bất định được dự liệu là giả thuyết có một phần cơ sở thực nghiệm.

{\footnotesize
\setlength{\tabcolsep}{3pt}\renewcommand{\arraystretch}{1.15}
\begin{longtable}{>{\raggedright}p{0.5cm}>{\raggedright}p{4.0cm}>{\raggedright}p{5.6cm}>{\raggedright}p{2.3cm}>{\raggedright\arraybackslash}p{1.7cm}}
\toprule
TT & Khẳng định & Thí nghiệm và điều đã đo & Nguồn & Trạng thái \\
\midrule\endfirsthead
\toprule
TT & Khẳng định & Thí nghiệm và điều đã đo & Nguồn & Trạng thái \\
\midrule\endhead
1 & Não có ít nơ-ron \nt{ACET}, chúng tập trung trong vài nhóm, còn đầu ra tỏa khắp vỏ não: một kênh quảng bá & Nhuộm kháng thể với enzyme tổng hợp acetylcholine và đánh dấu ngược dòng ở chuột cống: vỏ não, hồi hải mã, hạch hạnh nhân, hành khứu và đồi thị nhận acetylcholine chủ yếu từ sáu nhóm tế bào đặc của não trước nền và thân não & \cite{Mesulam1983} & đã đo \\
2 & Mức acetylcholine trong vỏ não cao khi thức và thấp trong giấc ngủ sóng chậm & Vi thẩm tách ở vỏ não và hồi hải mã của mèo di chuyển tự do, kèm ghi điện não đồ: lượng acetylcholine phóng ra khi thức yên tĩnh cao hơn $\sim100\%$, khi thức hoạt động cao hơn $\sim175\%$ so với giấc ngủ sóng chậm; trong giấc ngủ REM, ở vỏ não bằng mức khi thức & \cite{Marrosu1995,Hasselmo1999} & đã đo \\
3 & Ý nghĩa của tín hiệu do thụ thể quyết định: acetylcholine có các thụ thể--kênh nhanh và các thụ thể chậm khởi động phản ứng trong tế bào (mệnh đề~\ref{prop:receptor}) & Tổng quan các phép đo: tác động của acetylcholine lên tính dễ kích thích, truyền dẫn và tính dẻo phụ thuộc vào nơi phóng, phân nhóm thụ thể (nicotinic là kênh ion, muscarinic hoạt động qua protein G) và tế bào nhận & \cite{Picciotto2012} & đã đo \\
4 & Tác động thứ nhất: làm yếu liên kết nội bộ, gần như không động tới đầu vào từ bên ngoài (thừa số $1-a$ trong~\eqref{eq:acetnet}) & Lát cắt vỏ não khứu giác chuột cống: với $100$\,\textmu M chất chủ vận acetylcholine, điện thế của liên kết nội bộ (lớp Ib) giảm hơn $60\%$, của đầu vào bên ngoài (lớp Ia) --- dưới $15\%$, trong cùng một tế bào; sự ức chế là ở phía trước khớp thần kinh. Lát cắt hồi hải mã (CA1): $89\%$ so với $40\%$ cho đường nội bộ và đường vào & \cite{HasselmoBower1992,HasselmoSchnell1994} & đã đo \\
5 & Tác động thứ hai: tăng cường đầu vào (thừa số $\frac{1+a}{2}$) & Lát cắt tân vỏ não: thụ thể nicotinic tăng cường các khớp thần kinh của đầu vào từ đồi thị và không tác động lên các khớp thần kinh trong vỏ não; thụ thể muscarinic làm yếu cả hai loại. Acetylcholine làm tăng tính dễ kích thích của nơ-ron (tổng quan) & \cite{Gil1997,Hasselmo2006} & đã đo \\
6 & Tác động thứ ba: làm trọng số dễ thay đổi hơn (tốc độ $\eta a$) & Chuột cống: kích thích điện nhân nền (nguồn acetylcholine cho vỏ não) ghép cặp với một âm gây ra ở con vật trưởng thành sự tái tổ chức mạnh bản đồ vỏ não thính giác theo âm được trình bày & \cite{KilgardMerzenich1998,Hasselmo2006} & đã đo \\
7 & Quy tắc Hebb cho mẫu thưa trong phương trình thứ hai của~\eqref{eq:acetnet} cho ra một bộ nhớ liên kết, điền nốt mẫu từ một phần & Tính dung lượng của mạng với mẫu thưa và quy tắc hiệp phương sai: khi tỷ lệ nơ-ron hoạt động $p$ nhỏ, mạng lưu được nhiều mẫu hơn đáng kể so với khi $p=1/2$ & \cite{Tsodyks1988} & lý thuyết \\
8 & Ghi khi $a$ thấp làm hỏng bộ nhớ: mạng ghi điều nó nhớ lại, không phải điều nó thấy; khi $a=0.1$, hư hỏng không nhẹ hơn khi $a=0$ & \texttt{models/\allowbreak acet\_memory.py}: $200$ nơ-ron, năm mẫu mỗi mẫu $20$, mẫu mới có chung $10$ nơ-ron với một trong số đó, $10$ lần chạy. Khi $a=0$, mẫu mới không được ghi nhớ, mẫu cũ kéo theo $5.9$ trong $10$ nơ-ron lạ; khi $a=0.1$, mẫu mới được nhớ lại ($0.97$), nhưng hai mẫu hòa lẫn: mẫu mới kéo theo $7.3$ nơ-ron lạ, mẫu cũ --- $7.9$; từ $a=0.2$ --- ít hơn một & suy ra trong chương & mô hình \\
9 & Mệnh đề~\ref{prop:writeread}: không có mức $a$ nào mà cả ghi lẫn đọc đều tốt như nhau (hình~\ref{fig:acetsim}) & Cùng script, tốc độ học $\eta a$: mẫu mới được ghi nhớ trọn vẹn sau một lần trình bày khi $a\ge0.9$, được $0.8$ khi $a=0.75$, không được ghi nhớ khi $a\le0.5$. Đọc từ một nửa mẫu: $1.0$ khi $a\le0.2$, $0.96$ khi $a=0.3$, $0.04$ khi $a=0.4$ & suy ra trong chương & mô hình \\
10 & Trong não, một dây quảng bá --- \nt{ACET} --- chuyển giữa ghi và đọc & Ý tưởng lấy từ các mô hình xây dựng dựa trên phép đo trên lát cắt. Thí nghiệm trực tiếp nhất là ở người: chất chẹn thụ thể muscarinic scopolamine làm giảm khả năng học thuộc các cặp từ mới, mạnh nhất với những cặp chồng lấn với cặp cũ, nhưng không làm giảm việc nhớ lại các cặp đã học. Chưa có phép đo trực tiếp hai chế độ của mạng & \cite{HasselmoAndersonBower1992,Atri2004} & giả thuyết \\
11 & Bộ lọc~\eqref{eq:kalman} là tối ưu; trọng số đầu vào và tốc độ học là cùng một đại lượng $P/R$, ở chế độ xác lập $P=\sqrt{qR}$ và thời gian nhớ $\sqrt{R/q}$ (mệnh đề~\ref{prop:kalman}) & Không cần thí nghiệm: tính tối ưu của bộ lọc Kalman--Bucy là kết quả kinh điển của lý thuyết ước lượng; chế độ xác lập được suy ra trong chương & \cite{KalmanBucy1961} & lý thuyết \\
12 & \nt{ACET} báo cho mạng một đại lượng giống $P$ --- sự bất định được dự liệu & Được đề xuất trong một mô hình xác suất về sự chú ý. Chưa có phép đo trực tiếp cho thấy mức acetylcholine bám theo độ bất định của ước lượng & \cite{YuDayan2005} & giả thuyết \\
13 & Một phần nơ-ron \nt{ACET} đáp ứng với phần thưởng và hình phạt trong khoảng hai chục mili giây, càng mạnh khi sự củng cố càng bất ngờ & Chuột nhắt, nơ-ron cholinergic của não trước nền được nhận dạng bằng quang di truyền: đáp ứng với phần thưởng và hình phạt sau $18\pm3$\,ms, độ lớn tăng theo mức bất ngờ của sự củng cố, đáp ứng giống nhau ở nơ-ron của hai nhân & \cite{Hangya2015} & đã đo \\
14 & \nt{NORA} nhận ra thay đổi bất ngờ của thế giới, \nt{ACET} giữ mạng ở chế độ ghi & Sự phân công được đề xuất trong một mô hình. Gián tiếp: ở người, đường kính đồng tử, liên quan tới \nt{NORA}, bám theo các thay đổi và tốc độ học. Chưa có kiểm tra trực tiếp cặp \nt{NORA}--\nt{ACET} & \cite{YuDayan2005,Nassar2012} & giả thuyết \\
15 & Trong giấc ngủ sóng chậm, mạng ở chế độ đọc phát lại những gì được ghi ban ngày, và những lần phát lại này chép chúng sang bộ nhớ dài hạn & Ghi các tập hợp nơ-ron hồi hải mã ở chuột cống: các tế bào cùng hoạt động ban ngày thường cùng phát xung hơn trong giấc ngủ sóng chậm sau đó, và theo cùng thứ tự --- đã đo. Về việc chép: ức chế các sóng gợn hồi hải mã sau khi học làm giảm trí nhớ không gian, còn kéo dài các sóng gợn tự phát thì cải thiện nó; việc nguyên nhân là mức \nt{ACET} thấp chưa được chứng minh & \cite{WilsonMcNaughton1994,Skaggs1996,Girardeau2009,FernandezRuiz2019,Hasselmo1999} & giả thuyết \\
\bottomrule
\end{longtable}}
